\documentclass[11pt]{article}
\usepackage[margin=1in]{geometry}
\usepackage{amsmath,amssymb}
\usepackage{booktabs}
\usepackage{graphicx}
\usepackage{microtype}
\usepackage[hidelinks]{hyperref}
\usepackage{enumitem}
\usepackage{caption}
\usepackage{subcaption}
\usepackage{array}
\usepackage{multirow}
\usepackage{xcolor}
\usepackage{titlesec}
\usepackage[numbers,sort&compress]{natbib}
\usepackage{fancyhdr}

\captionsetup{font=small,labelfont=bf}
\setlist[itemize]{leftmargin=1.2em,itemsep=2pt,topsep=3pt}
\setlist[enumerate]{leftmargin=1.4em,itemsep=2pt,topsep=3pt}
\titleformat{\section}{\large\bfseries}{\thesection}{0.6em}{}
\titleformat{\subsection}{\normalsize\bfseries}{\thesubsection}{0.6em}{}
\renewcommand{\arraystretch}{1.08}
\pagestyle{fancy}
\fancyhf{}
\fancyhead[L]{\footnotesize Verification Debt}
\fancyhead[R]{\footnotesize\thepage}
\renewcommand{\headrulewidth}{0.3pt}

% generated by emit_numbers.py - do not edit by hand
% source: results/summary.json  (480 records)

\newcommand{\VDnumRecords}{480}
\newcommand{\VDnBase}{35}
\newcommand{\VDnGuard}{25}
\newcommand{\VDnFamilies}{4}
\newcommand{\VDrhoIn}{3.23}
\newcommand{\VDrhoOut}{3.25}
\newcommand{\VDsmin}{774}
\newcommand{\VDsmax}{923}
\newcommand{\VDsratio}{1.19}
\newcommand{\VDspearman}{0.40}
\newcommand{\VDpperm}{0.75}
\newcommand{\VDbarB}{met}
\newcommand{\VDbarC}{missed}
\newcommand{\VDbarD}{missed}
\newcommand{\VDworstFamily}{Invoices}
\newcommand{\VDWorstD}{8.89}
\newcommand{\VDWorstDlo}{7.83}
\newcommand{\VDWorstDhi}{9.97}
\newcommand{\VDWorstGuardCut}{-112\%}
\newcommand{\VDWorstGuardCost}{101\%}
\newcommand{\VDWorstGuardD}{18.1}
\newcommand{\VDWorstGuardBase}{8.56}
\newcommand{\VDWorstGuardDlo}{7.30}
\newcommand{\VDWorstGuardDhi}{11.9}
\newcommand{\VDCleanseS}{56.0}
\newcommand{\VDCleanseV}{187.8}
\newcommand{\VDCleanseR}{195.1}
\newcommand{\VDCleanseD}{6.84}
\newcommand{\VDCleanseDlo}{6.09}
\newcommand{\VDCleanseDhi}{7.69}
\newcommand{\VDCleanseDmean}{6.83}
\newcommand{\VDCleanseVS}{3.35}
\newcommand{\VDCleanseRS}{3.48}
\newcommand{\VDCleanseVis}{0.778}
\newcommand{\VDCleansePass}{31\%}
\newcommand{\VDCleanseFinal}{41\%}
\newcommand{\VDCleanseVerAcc}{0\%}
\newcommand{\VDCleanseRwRate}{100\%}
\newcommand{\VDCleanseRounds}{1.50}
\newcommand{\VDCleanseEffort}{6651}
\newcommand{\VDCleanseDfull}{8.31}
\newcommand{\VDCleanseN}{70}
\newcommand{\VDEmailsS}{61.4}
\newcommand{\VDEmailsV}{233.6}
\newcommand{\VDEmailsR}{172.7}
\newcommand{\VDEmailsD}{6.62}
\newcommand{\VDEmailsDlo}{5.95}
\newcommand{\VDEmailsDhi}{7.34}
\newcommand{\VDEmailsDmean}{6.62}
\newcommand{\VDEmailsVS}{3.81}
\newcommand{\VDEmailsRS}{2.81}
\newcommand{\VDEmailsVis}{0.556}
\newcommand{\VDEmailsPass}{0\%}
\newcommand{\VDEmailsFinal}{0\%}
\newcommand{\VDEmailsVerAcc}{1\%}
\newcommand{\VDEmailsRwRate}{99\%}
\newcommand{\VDEmailsRounds}{1.71}
\newcommand{\VDEmailsEffort}{6776}
\newcommand{\VDEmailsDfull}{7.73}
\newcommand{\VDEmailsN}{70}
\newcommand{\VDInvoicesS}{64.6}
\newcommand{\VDInvoicesV}{299.1}
\newcommand{\VDInvoicesR}{275.4}
\newcommand{\VDInvoicesD}{8.89}
\newcommand{\VDInvoicesDlo}{7.83}
\newcommand{\VDInvoicesDhi}{9.97}
\newcommand{\VDInvoicesDmean}{8.89}
\newcommand{\VDInvoicesVS}{4.63}
\newcommand{\VDInvoicesRS}{4.26}
\newcommand{\VDInvoicesVis}{1.000}
\newcommand{\VDInvoicesPass}{0\%}
\newcommand{\VDInvoicesFinal}{0\%}
\newcommand{\VDInvoicesVerAcc}{0\%}
\newcommand{\VDInvoicesRwRate}{100\%}
\newcommand{\VDInvoicesRounds}{2.39}
\newcommand{\VDInvoicesEffort}{9339}
\newcommand{\VDInvoicesDfull}{10.1}
\newcommand{\VDInvoicesN}{70}
\newcommand{\VDReportsS}{54.1}
\newcommand{\VDReportsV}{233.1}
\newcommand{\VDReportsR}{193.7}
\newcommand{\VDReportsD}{7.88}
\newcommand{\VDReportsDlo}{6.90}
\newcommand{\VDReportsDhi}{8.96}
\newcommand{\VDReportsDmean}{7.88}
\newcommand{\VDReportsVS}{4.30}
\newcommand{\VDReportsRS}{3.58}
\newcommand{\VDReportsVis}{0.222}
\newcommand{\VDReportsPass}{0\%}
\newcommand{\VDReportsFinal}{0\%}
\newcommand{\VDReportsVerAcc}{1\%}
\newcommand{\VDReportsRwRate}{99\%}
\newcommand{\VDReportsRounds}{1.43}
\newcommand{\VDReportsEffort}{7525}
\newcommand{\VDReportsDfull}{9.73}
\newcommand{\VDReportsN}{70}
\newcommand{\VDDmin}{6.62}
\newcommand{\VDDmax}{8.89}
\newcommand{\VDdfullMin}{7.73}
\newcommand{\VDdfullMax}{10.1}
\newcommand{\VDfirstPass}{22}
\newcommand{\VDfirstPassRate}{8\%}
\newcommand{\VDCleanseGuardBase}{6.82}
\newcommand{\VDCleanseGuardD}{11.7}
\newcommand{\VDCleanseGuardDelta}{-72\%}
\newcommand{\VDCleanseGuardChange}{72\%}
\newcommand{\VDCleanseGuardCost}{63\%}
\newcommand{\VDCleanseGuardPass}{40\%}
\newcommand{\VDCleanseGuardDlo}{3.06}
\newcommand{\VDCleanseGuardDhi}{6.73}
\newcommand{\VDEmailsGuardBase}{6.50}
\newcommand{\VDEmailsGuardD}{11.3}
\newcommand{\VDEmailsGuardDelta}{-74\%}
\newcommand{\VDEmailsGuardChange}{74\%}
\newcommand{\VDEmailsGuardCost}{68\%}
\newcommand{\VDEmailsGuardPass}{0\%}
\newcommand{\VDEmailsGuardDlo}{3.22}
\newcommand{\VDEmailsGuardDhi}{6.46}
\newcommand{\VDInvoicesGuardBase}{8.56}
\newcommand{\VDInvoicesGuardD}{18.1}
\newcommand{\VDInvoicesGuardDelta}{-112\%}
\newcommand{\VDInvoicesGuardChange}{112\%}
\newcommand{\VDInvoicesGuardCost}{101\%}
\newcommand{\VDInvoicesGuardPass}{0\%}
\newcommand{\VDInvoicesGuardDlo}{7.30}
\newcommand{\VDInvoicesGuardDhi}{11.9}
\newcommand{\VDReportsGuardBase}{8.08}
\newcommand{\VDReportsGuardD}{15.8}
\newcommand{\VDReportsGuardDelta}{-95\%}
\newcommand{\VDReportsGuardChange}{95\%}
\newcommand{\VDReportsGuardCost}{81\%}
\newcommand{\VDReportsGuardPass}{0\%}
\newcommand{\VDReportsGuardDlo}{5.39}
\newcommand{\VDReportsGuardDhi}{10.1}
\newcommand{\VDguardDeltaLo}{-112\%}
\newcommand{\VDguardDeltaHi}{-72\%}
\newcommand{\VDguardRiseLo}{72\%}
\newcommand{\VDguardRiseHi}{112\%}
\newcommand{\VDguardCostLo}{63\%}
\newcommand{\VDguardCostHi}{101\%}
\newcommand{\VDqDCleanse}{7.00}
\newcommand{\VDqDEmails}{7.49}
\newcommand{\VDqDInvoices}{11.7}
\newcommand{\VDqDReports}{9.62}
\newcommand{\VDqSpearman}{0.20}
\newcommand{\VDqPperm}{0.92}
\newcommand{\VDlDCleanse}{6.68}
\newcommand{\VDlDEmails}{5.74}
\newcommand{\VDlDInvoices}{6.06}
\newcommand{\VDlDReports}{6.15}
\newcommand{\VDlSpearman}{0.00}
\newcommand{\VDlPperm}{1.00}
\newcommand{\VDqLeak}{100.0\%}
\newcommand{\VDqLeakCrit}{100.0\%}
\newcommand{\VDqLeakN}{2}
\newcommand{\VDlLeak}{n/a}
\newcommand{\VDlLeakCrit}{n/a}
\newcommand{\VDlLeakN}{0}
\newcommand{\VDsensMin}{5.86}
\newcommand{\VDsensMax}{46.9}
\newcommand{\VDtermCapA}{4.60}
\newcommand{\VDtermCapB}{6.80}
\newcommand{\VDtermCapC}{7.87}
\newcommand{\VDtermCapE}{8.89}
\newcommand{\VDtermPatA}{5.55}
\newcommand{\VDdetTP}{256}
\newcommand{\VDdetFP}{22}
\newcommand{\VDdetFN}{2}
\newcommand{\VDdetTN}{0}
\newcommand{\VDdetPrec}{92\%}
\newcommand{\VDdetRec}{99\%}
\newcommand{\VDdetFalseRej}{8\%}
\newcommand{\VDdetReworkFalse}{6\%}
\newcommand{\VDdetCleansePrec}{69\%}
\newcommand{\VDdetCleanseFalseRej}{31\%}
\newcommand{\VDdetCleanseReworkFalse}{27\%}
\newcommand{\VDdetEmailsPrec}{100\%}
\newcommand{\VDdetEmailsFalseRej}{0\%}
\newcommand{\VDdetEmailsReworkFalse}{0\%}
\newcommand{\VDdetInvoicesPrec}{100\%}
\newcommand{\VDdetInvoicesFalseRej}{0\%}
\newcommand{\VDdetInvoicesReworkFalse}{0\%}
\newcommand{\VDdetReportsPrec}{100\%}
\newcommand{\VDdetReportsFalseRej}{0\%}
\newcommand{\VDdetReportsReworkFalse}{0\%}
\newcommand{\VDnBaseline}{280}
\newcommand{\VDBaseAccPooled}{2/280}
\newcommand{\VDchkShare}{6\%}
\newcommand{\VDchkShareLo}{3\%}
\newcommand{\VDchkShareHi}{9\%}
\newcommand{\VDchkD}{0.45}
\newcommand{\VDchkDLo}{0.27}
\newcommand{\VDchkDHi}{0.65}
\newcommand{\VDchkDTotal}{7.58}
\newcommand{\VDchecksCleanse}{30}
\newcommand{\VDchecksEmails}{12}
\newcommand{\VDchecksInvoices}{18}
\newcommand{\VDchecksReports}{10}
\newcommand{\VDlambdaStarMin}{0.11}
\newcommand{\VDlambdaStarMax}{0.15}
\newcommand{\VDlambdaStarQMin}{0.085}
\newcommand{\VDlambdaStarQMax}{0.143}
\newcommand{\VDpromptShareMin}{80\%}
\newcommand{\VDpromptShareMax}{89\%}
\newcommand{\VDverRecall}{99.2\%}
\newcommand{\VDverCaught}{256}
\newcommand{\VDverBlocked}{258}
\newcommand{\VDgoldDmin}{1.44}
\newcommand{\VDgoldDmax}{2.25}
\newcommand{\VDgoldFrMin}{80\%}
\newcommand{\VDgoldFrMax}{100\%}
\newcommand{\VDgoldRej}{39/40}
\newcommand{\VDgoldAcc}{1/40}
\newcommand{\VDgoldRejPct}{98\%}
\newcommand{\VDgoldRejRate}{97.5\%}
\newcommand{\VDgoldRejLo}{87\%}
\newcommand{\VDgoldRejHi}{100\%}
\newcommand{\VDGoldDllamaCleanse}{1.83}
\newcommand{\VDGoldFrllamaCleanse}{100\%}
\newcommand{\VDGoldAccllamaCleanse}{0/5}
\newcommand{\VDGoldDllamaEmails}{1.48}
\newcommand{\VDGoldFrllamaEmails}{100\%}
\newcommand{\VDGoldAccllamaEmails}{0/5}
\newcommand{\VDGoldDllamaInvoices}{1.78}
\newcommand{\VDGoldFrllamaInvoices}{100\%}
\newcommand{\VDGoldAccllamaInvoices}{0/5}
\newcommand{\VDGoldDllamaReports}{2.20}
\newcommand{\VDGoldFrllamaReports}{100\%}
\newcommand{\VDGoldAccllamaReports}{0/5}
\newcommand{\VDGoldDqwenCleanse}{1.97}
\newcommand{\VDGoldFrqwenCleanse}{100\%}
\newcommand{\VDGoldAccqwenCleanse}{0/5}
\newcommand{\VDGoldDqwenEmails}{1.44}
\newcommand{\VDGoldFrqwenEmails}{80\%}
\newcommand{\VDGoldAccqwenEmails}{1/5}
\newcommand{\VDGoldDqwenInvoices}{1.88}
\newcommand{\VDGoldFrqwenInvoices}{100\%}
\newcommand{\VDGoldAccqwenInvoices}{0/5}
\newcommand{\VDGoldDqwenReports}{2.25}
\newcommand{\VDGoldFrqwenReports}{100\%}
\newcommand{\VDGoldAccqwenReports}{0/5}
\newcommand{\VDGoldDCleanse}{1.97}
\newcommand{\VDGoldFrCleanse}{100\%}
\newcommand{\VDGoldAccCleanse}{0/5}
\newcommand{\VDGoldDEmails}{1.44}
\newcommand{\VDGoldFrEmails}{80\%}
\newcommand{\VDGoldAccEmails}{1/5}
\newcommand{\VDGoldDInvoices}{1.88}
\newcommand{\VDGoldFrInvoices}{100\%}
\newcommand{\VDGoldAccInvoices}{0/5}
\newcommand{\VDGoldDReports}{2.25}
\newcommand{\VDGoldFrReports}{100\%}
\newcommand{\VDGoldAccReports}{0/5}
\newcommand{\VDqBaseAcc}{2/140}
\newcommand{\VDlBaseAcc}{0/140}
\newcommand{\VDgoldFisher}{0.33}


\newcommand{\D}{\ensuremath{D}}
\newcommand{\SV}{\ensuremath{S}}
\newcommand{\VV}{\ensuremath{V}}
\newcommand{\RV}{\ensuremath{R}}
\newcommand{\famname}[1]{\textsc{#1}}

\title{\bfseries Verification Debt:\\[2pt]
The Hidden Balance Sheet of Agent-Produced Work}
\author{Boyu Wang\\[2pt]
{\small School of Management, University of San Francisco, United States,
\texttt{bwang62@usfca.edu}}}
\date{October 2026}

\begin{document}
\maketitle

\begin{abstract}
\noindent
Agent deployment does not eliminate the cost of work; it moves it downstream, into checking what
the agent produced and repairing what fails. We give it an accounting construct. For a task family,
$\SV$ is the production saving the agent displaces against a stated baseline, $\VV$ is the effort
to check a product to a fixed acceptance standard, and $\RV$ is the effort to repair failures. The
debt ratio is $\D=(\VV+\RV)/\SV$, and $\SV-(\VV+\RV)=\SV(1-\D)$ shows what production-only ROI
leaves out. We apply a protocol --- planted ground truth, a deterministic defect census, an
independent verification pass, and rework loops --- to four synthetic process families on two local
open-weight models (\VDnumRecords{} work products, \VDnBaseline{} in the baseline arm). First-pass
acceptance against the census was rare: \VDfirstPass{} of \VDnBaseline{} baseline products passed
(concentrated in \famname{cleansing}),
and the verifier accepted essentially none of even the clean ones, so $\D$ is a multiple of the
saving on
every family (\VDDmin--\VDDmax). A positive control reframes this: the planted gold deliverable is
census-perfect, yet the verifier rejects \VDgoldRej{} of them (\VDgoldRejRate{}; 95\% CI
\VDgoldRejLo--\VDgoldRejHi). A
single verification pass already gives $\D_{\text{gold}}=\VDgoldDmin$--$\VDgoldDmax$, so the excess
of the measured $\D$ over $\D_{\text{gold}}$ includes verifier over-rejection on clean work, which
we quantify: \VDdetFalseRej{} of first-pass rejections are false, concentrated on \famname{cleansing}
and gold, and the excess over a one-pass gold baseline (\S\ref{sec:proxy}) of trajectories with a
false first rejection is \VDchkShare{} of all verification-plus-rework tokens (95\% CI
\VDchkShareLo--\VDchkShareHi) --- so checker error inflates $\D$ but is a minority of it. $\D\geq 1$
is therefore guaranteed by the protocol's unit asymmetry (Section~\ref{sec:proxy}), not by
the work. The verifiability gradient was
absent --- no detectable monotonic gradient ($\rho_s=\VDspearman$), and a guardrail increased $\D$ by
\VDguardRiseLo--\VDguardRiseHi. Code, logs and the regeneration script are released alongside
this paper (Appendix~\ref{app:repro}).
\end{abstract}

\section{Introduction}

A firm that deploys an agent to process invoices, draft support replies, clean records, or
summarise operations observes the same thing at first: work that used to take an afternoon
takes minutes. What it observes next is less comfortable. The invoices are mostly right, but
someone has to find the ones that are not. The support drafts are fluent, but a draft that
promises a refund the policy forbids costs more than a slow, correct one. The summary is
well written and contains a number that appears nowhere in the table.

None of that cost is in the production metric. The production metric counts the work that was
produced. The cost of establishing whether the work is any good, and of repairing it when it is
not, lands on a different line, later, and usually on a different team. When the two lines are
not added together, the reported gain is not a gain. It is a partial sum.

This paper is about adding the lines together, for a unit of agent-produced work, in a way that
can be measured rather than argued about.

\paragraph{The claim, in one sentence.}
Agent deployment does not eliminate the cost of work --- it shifts cost downstream into
verification and rework, and on every task family we measured that shifted cost is a multiple of
the production saving. The more useful finding is what the multiple is made of: a structural floor
(the checker re-reads the inputs), plus a checker-error component that we now measure --- false
first-pass rejections are \VDdetFalseRej{} of rejections, and the excess over a one-pass gold
baseline (\S\ref{sec:proxy}) of trajectories with a false first rejection is \VDchkShare{} of all
$\VV+\RV$ (95\% CI \VDchkShareLo--\VDchkShareHi) --- and the rest is repair of genuine defects.

\paragraph{Why an accounting construct.}
The field already knows the symptom and the bottleneck. Firm-level evidence on coding agents
finds output up and delivered work flat, with review absorbing the difference
\citep{chen2026}. Case studies name review as the constraint. What is missing is not another
observation but a \emph{unit of account}: a quantity defined per unit of agent-produced work,
decomposed into its parts, measurable under controlled conditions, and comparable across task
types. Without such a quantity, every claim about agent ROI is a claim about a numerator with an
unstated denominator.

The construct we propose is deliberately simple. Three costs, one ratio. Its interest is not in
its arithmetic but in three properties: it is defined against a stated baseline rather than
against a vendor's counterfactual; it separates checking from repairing, because the two behave
differently under intervention; and it is parameterised by a property of the task, verifiability,
that can be rated before any model is run.

\paragraph{Contributions.}
\begin{enumerate}
\item \textbf{The construct and a structural bar.} We define, for a task family, the gross
production saving $\SV$, the verification cost $\VV$, the rework cost $\RV$, and the debt ratio
$\D=(\VV+\RV)/\SV$, together with the accounting identity that separates booked from net gain
(Section~\ref{sec:construct}). We observe $\D\geq 1$ on every family we measured, but report it as
a \emph{structural} consequence of the identity --- a perfect product still incurs the
verification-cost floor --- rather than as evidence about agent work quality
(Section~\ref{sec:results}).
\item \textbf{A pre-registered verifiability test and its (null) result.} We rate verifiability as
an independently rated property of the family before any model is run, state the ordering
hypothesis up front, and test it against measured ratios. The test \emph{failed}: the rank
correlation was absent, so verifiability did not predict where the debt landed
(Sections~\ref{sec:construct} and~\ref{sec:results}).
\item \textbf{The measurement protocol.} Planted ground truth, a deterministic defect census by
severity, an independent verification pass with a fixed rubric, and rework loops run to
acceptance or a round cap. The protocol, not only the numbers, is the deliverable: it transfers
to task families we did not build (Section~\ref{sec:protocol}).
\item \textbf{The intervention result.} Guardrails --- an acceptance checklist plus a mandated
dual verification pass --- treated as a debt-management instrument: we measure both the change in
$\D$ and the throughput it costs, and report that it moved $\D$ the wrong way on every family
(Section~\ref{sec:results}).
\item \textbf{A measured proxy and its bias.} Verification effort is measured in agent tokens.
We say what that proxy can and cannot support, give the direction of its bias, and report the
sensitivity band over conversion factors instead of silently converting agent time into human
hours (Section~\ref{sec:proxy}).
\end{enumerate}

\paragraph{What this paper does not claim.}
It does not measure any organisation's ROI. It does not claim that the laboratory ratios are the
field ratios. It does not claim that agents are unproductive. It does not claim that the magnitude
of the measured quantity reflects agent work quality --- $\D\geq 1$ is a structural consequence of
the protocol's unit asymmetry, not a finding about the work. And it does not claim that the verifiability
axis explains where the debt lands --- the pre-registered test found no detectable gradient, so it did
not.

\section{What the field sees, and what it is missing}
\label{sec:field}

\paragraph{The productivity paradox, again.}
The pattern in which measured output rises while measured organisational performance does not is
old enough to have a literature. \citet{brynjolfsson1993} named the productivity paradox --- large
IT investments coexisting with weak aggregate productivity gains, a discrepancy he attributed
partly to measurement; \citet{brynjolfsson1996} showed the firm-level returns were there,
conditional on complementary investment; and \citet{brynjolfsson2021} generalised the mechanism
into a J-curve, in which intangible complements are expensed before they pay back. The lesson is
not that the gains are illusory but that the accounting boundary determines the sign of the
answer. Agent deployment reproduces the structure with one difference: the omitted complement is
now \emph{labour spent checking}, diffuse and absorbed into existing roles, and so invisible in
exactly the way that makes an ROI deck look good.

\paragraph{The empirical hook.}
The clearest firm-level evidence to date is the staggered difference-in-differences study of
coding agents by \citet{chen2026}, ``Artificial Intelligence in the Firm: Bottlenecks in Software
Production,'' using engineering-management telemetry from 718 firms (more than 700{,}000 workers,
some 300 million work events).\footnote{Figures are as reported in the August 2026 version of the
working paper; we cite them as reported rather than re-derived from the microdata, which is not
public.} Adoption of coding agents is associated with a 30\% increase in lines of code, 20\% more
commits, and 23\% more pull requests, all statistically significant; autocomplete-only adoption
shows a smaller 12\% increase. Measures of delivered work show no statistically detectable change:
issue and epic resolution rates do not move, and neither does headcount. The interval rules out
growth above roughly 12\%, not small growth, so we read this as ``no detectable change'' rather
than ``flat.'' The difference is absorbed downstream. Pull-request time-to-merge rises 49\%, the
share of pull requests receiving change requests rises by about 12 percentage points (from a base
near 13\%), comments per pull request rise 35\%, and the share of engineers doing review rises
14\%. Field experiments complement the firm-level view from the individual side: three randomised
rollouts of GitHub Copilot across 4{,}867 developers at Microsoft, Accenture and an anonymous
Fortune~100 firm measured individual productivity effects directly \citep{cui2024}. Related field and survey evidence points the same way
\citep{demirer2026,metr2025,dora2025}. Both literatures
measure production; neither books
the checking bill, which is the measurement this paper attempts under controlled conditions.

Read as a production story, this is a puzzle. Read as an accounting story, it is not a puzzle at
all: production went up, the verification stage became the binding constraint, and the two effects
cancelled in the only number that was being counted. What the study cannot tell us is the
\emph{rate} at which the two effects cancel, because it observes neither the verification effort
nor the rework in comparable units, and because it observes one task family --- code --- whose
verifiability is unusually high. That is the gap this paper fills.

\paragraph{The gap, stated exactly.}
Field studies observe the symptom (output up, throughput flat) and case studies name the
bottleneck (review), but nobody has built the accounting construct --- a verification debt ratio
defined per unit of agent-produced work, decomposed into verification effort and rework, and
measured it across process-task families under controlled conditions with planted ground truth.
That is this paper.

\paragraph{Relation to workflow portfolios.}
A companion line of work treats deployment as a portfolio problem: which of several agentic
workflows to run, how many executions to buy, and how to select among their outputs when the
selector is imperfect \citep{abdolmaleki2026}. That work answers \emph{which workflow to run} and
prices selection error against compute. This paper is its accounting-side companion: it answers
\emph{what the run actually costs} once verification is on the books. The two are complementary in
a specific way. Portfolio selection chooses among workflows whose verification structures differ;
a debt ratio is what makes those differences comparable in money rather than in accuracy.

\paragraph{Relation to prior uses of the term.}
``Verification debt'' is already in circulation, and we do not claim the phrase. Industry writing
and practitioner commentary use it as a \emph{stock}: the accumulated
backlog of AI-generated code that has not been verified to the organisation's standard.
The phrase is Vogels's \citep{vogels2025}, and \citet{sonar2026} quantify the stock in a survey of more than
1{,}100 developers: 96\% do not fully trust AI-generated code, yet only 48\% say they always
verify it before committing. \citet{agilev2026} uses it similarly, as the shortfall that accrues when output grows faster than
verification capacity. Our quantity is a \emph{flow ratio} --- verification plus rework cost per
unit of production saving, with a break-even at one --- which is a different object, and the
distinction is the point: a stock tells you how much unverified work exists, a flow ratio tells you
whether a deployment creates value per unit of work. \citet{bhati2026} proposes a structurally
similar ratio (a ``verification tax'' with verification and rework in the numerator) but normalises
by generation cost rather than by the displaced production, which is what yields the net-gain
identity and the $D=1$ break-even here. \citet{crescitelli2026} argue that evaluation should measure
verification cost rather than correctness, and note that a ratio needs a common unit; we take that
objection seriously in Section~\ref{sec:proxy}.

\paragraph{Relation to the economics of verification.}
\citet{catalini2026} argue, at the level of the whole economy, that the binding constraint on
agent deployment is the cost to verify rather than the cost to automate, and that the gap
between the two is widening. Theirs is a model of the aggregate; it takes the cost to verify as
a primitive. This paper is the micro counterpart: an accounting identity and a measurement
protocol that price verification and rework per task family, and a gold control that separates
the part of the ratio any checking process must pay from the part a given deployment adds.
Where \citet{catalini2026} predict that verification binds, we measure by how much, and where
the floor sits.

\paragraph{Classical antecedents.}
The decomposition is the cost-of-quality framework with different names: $\VV$ is appraisal cost,
$\RV$ is internal-failure cost, and leakage is external-failure cost, the classical four-category
scheme of Juran and Feigenbaum \citep{juran1999,feigenbaum1961}. The verification-and-rework trade
also has a decision-theoretic ancestor in costly state verification \citep{townsend1979}, where a
principal pays to observe a state and the optimal audit is costly and imperfect. What is new here is
not the categories but the \emph{denominator}: pricing the burden against the production the agent
displaces turns the cost-of-quality accounting into a per-unit break-even that can be compared
across task families and checked against a pre-registered ordering. We say this plainly because a
referee will, and because the contribution is stronger when stated as a measurement protocol built
on old categories than as a new concept.

\section{The construct}
\label{sec:construct}

\subsection{Three costs and one ratio}

Fix a task family $f$ --- a class of work products that a procedure can produce from a stated
input. For a deployment configuration $c$ (a model, a prompt, a decoding setting, a workflow), and
for a unit of agent-produced work, define:

\begin{description}[leftmargin=1.4em,style=nextline,itemsep=3pt]
\item[Gross production saving $\SV$.] The human-equivalent production cost that the agent
displaces, measured against a \emph{stated} baseline procedure: the effort a competent operator
would spend reading the raw inputs and writing the complete correct deliverable. $\SV$ is a
property of the task instance and the baseline procedure, not of the model, so the same
denominator applies to every arm compared on that instance.
\item[Verification cost $\VV$.] The effort required to check the product to a fixed acceptance
standard --- not to a standard the checker chooses after seeing the product. The standard is
written down before the run; the checker applies it.
\item[Rework cost $\RV$.] The effort required to bring failed products to acceptance, including
the regenerated volume. A product that fails three times and is then accepted has consumed three
repair attempts, and all three are cost.
\end{description}

The debt ratio is
\begin{equation}
\D \;=\; \frac{\VV+\RV}{\SV}.
\label{eq:debt}
\end{equation}
Because booked ROI counts $\SV$ and the deployment actually nets $\SV-(\VV+\RV)$, we have the
identity
\begin{equation}
\text{net gain} \;=\; \SV\,(1-\D).
\label{eq:identity}
\end{equation}
Equation~\eqref{eq:identity} is the whole argument in one line. If $\D=0.3$, a booked saving of
one million overstates the net gain by 43\%. If $\D=1$, the deployment is value-neutral and the
booked saving is entirely fictitious. If $\D>1$, the deployment destroys value on that family
under that configuration, and no amount of enthusiasm about the production number changes the
sign. None of this requires a claim about any particular firm: it is an identity, and the
empirical question is only the value of $\D$.

Two features of the definition are worth defending.

First, $\SV$ is defined against a \emph{procedure}, not against a wage bill. The alternative ---
pricing the saving as ``hours saved times loaded cost'' --- imports a counterfactual that nobody
measures: the hours a person would have spent had the agent not existed. Our baseline is
narrower and auditable: the volume of work the artifact contains, priced in the same unit as the
costs. Section~\ref{sec:proxy} states exactly what that unit is and what it cannot do.

Second, $\VV$ and $\RV$ are separated even though both are labour. They respond differently to
intervention. An acceptance checklist attacks $\RV$ by raising first-pass quality; a second
verification pass attacks $\RV$ by catching more before delivery but raises $\VV$ by construction.
A single combined number would hide the trade, which is the one a manager actually faces.

\subsection{Verifiability, and the ordering hypothesis}

Not all work is equally checkable. An invoice either matches the purchase order or it does not,
and a clerk with a spreadsheet settles it in seconds. A written summary of operational data may
contain a number that is wrong in a way that reads perfectly well, and settling it requires
re-deriving the number from the table. We define \emph{task verifiability} as an independently
rated property of the family --- the degree to which a checker can confirm correctness cheaply and
objectively --- and we operationalise it with a three-item rubric scored before any model is run
(Section~\ref{sec:protocol}). One scope restriction matters and is easy to miss: the rubric scores
checkability against a \emph{deterministic oracle} that knows the planted truth. Checkability for a
fallible language-model verifier is a different quantity, and the positive control
(Section~\ref{sec:results}) shows the two can diverge sharply --- work that is perfectly checkable
in principle can still be mis-graded by the checker actually deployed.

The hypothesis, stated before the experiment: \textbf{$\D$ increases as verifiability decreases.}
The mechanism has two channels, and they point the same way.
\begin{itemize}
\item $\VV$ rises, because checking low-verifiability work is nearly as hard as doing it. The
checker cannot compare against a key; it must re-derive.
\item $\RV$ rises, because defects in low-verifiability work are subtler. They survive the first
pass, so more of them reach rework, and the rework is harder to target because the defect is
often a judgement rather than a value.
\end{itemize}
The hypothesis is falsifiable in a useful way: it predicts a monotone ordering across families
with different verifiability, not merely a positive average.

\subsection{What $\D$ is not}

$\D$ is not a measure of model quality, and it is not comparable across studies that use different
baseline procedures. It is defined relative to a configuration: the same family can have
$\D<1$ with a strong model and a deterministic checker and $\D>1$ with a weak model and a
judgement-based checker. That is a feature. The quantity a manager needs is not ``is this task
worth automating'' in the abstract but ``is this task worth automating \emph{with this workflow
and this checking process}'', and $\D$ is indexed exactly that way.

Finally, $\D$ as defined excludes the compute cost of production itself. It is therefore a
\emph{lower bound} on the total displacement cost. We report the extended ratio
$\D^{\text{full}}=(P+\VV+\RV)/\SV$, where $P$ is the production cost, as an exploratory
variant in Section~\ref{sec:robust}.

\section{Measurement protocol}
\label{sec:protocol}

The protocol is the part of this paper intended to outlive its numbers. It has five components,
and each exists to close a specific way of fooling oneself.

\subsection{Planted ground truth}

Every instance is generated together with its correct deliverable. The generator is the only
source of truth: nothing is inferred from model output, and no human judges whether a product is
right. This closes the most common failure mode in agent evaluation, in which the grader is
another model whose errors are correlated with the producer's.

Four families were built, each a recognisable class of process work with defects planted at known
locations: invoice three-way match (18 graded fields), policy-constrained email resolutions (12),
record consolidation into entities (30), and operational report summaries (10). Appendix
Table~\ref{tab:spec} gives the full specification. The point for the protocol is only that each
product is a structured object the census can grade field by field, so a defect is a comparison
rather than a judgement.

\subsection{The defect census}

The census grades a product against planted truth with deterministic Python. It never sees a
prompt, a reasoning trace, or a token count, and it is implemented in a module that does not import
the producer's code. Defects are counted by severity under rules fixed per family before the run:
in every family, approving defective work is \emph{critical} and rejecting sound work is
\emph{major}, so that the census prices the error the organisation actually pays for rather than
the error that is easiest to count. Acceptance on the census side means no critical and no major
defect.

The census is validated before use: on the canonical (gold) deliverable it must return zero
defects, and on a deliberately corrupted deliverable it must return at least one blocking defect.
Both checks run in \texttt{scripts/selftest.py} and pass for all four families. Without that
check, a ground-truth bug and a model error are indistinguishable.

\subsection{The independent verification pass}

Verification is a separate pass with a different system prompt, no access to the producer's
reasoning, and an explicit mandate to re-derive the answer from the raw inputs before comparing.
It applies the family's fixed rubric and emits a verdict plus a list of defects with field,
severity, and its own recomputed value.

The pass is \emph{not} told the ground truth, and its verdict --- not the census --- gates the
rework loop, because that is what a deployment does: the organisation repairs what its checking
process catches. The gap between what the verifier accepts and what the census accepts is then
measured rather than assumed, and reported as \emph{leakage}
(Section~\ref{sec:leakage}). A weak checker does not make debt disappear; it moves debt off the
books, which is worse.

\subsection{Rework to acceptance}

A rejected product returns to the producer with the reviewer's findings and an instruction to
recompute independently --- where the producer's own arithmetic disagrees with the reviewer's
suggested value, it keeps its own. The repaired product is re-verified, and the loop continues
until acceptance or a cap of five rounds.

Two loop-termination rules are pre-registered alongside the cap. The loop stops early if a round
returns a byte-identical product, or if two consecutive rounds fail to reduce the number of
distinct findings. Both model what a deployment does when a repair loop stops making progress:
it escalates to a human. Recording the stall is part of the measurement, not a workaround --- a
debt that cannot be serviced by rework is a different object from a debt that is merely expensive
to service.

\subsection{Effort units}

All three components are expressed in one unit: \emph{tokens of text processed}, the token volume a
procedure must read plus write to do its job. $\VV$ and $\RV$ are measured exactly by the local
model server (prompt plus completion tokens, summed over every pass and every repair). $\SV$ is
computed from character volume and converted with a characters-per-token ratio calibrated on the
run's own traffic, separately for reading ($\rho_{\text{in}}=\VDrhoIn$) and writing
($\rho_{\text{out}}=\VDrhoOut$):
\begin{equation}
\SV \;=\; \frac{\lvert\text{raw inputs}\rvert}{\rho_{\text{in}}}
\;+\; \frac{\lvert\text{canonical deliverable}\rvert}{\rho_{\text{out}}}.
\label{eq:S}
\end{equation}
The canonical deliverable is the complete correct product in the family's output schema, including
a one-line justification per field, produced by the generator. $\SV$ therefore prices the volume of
work the artifact contains. Wall-clock latency is logged per record, but only in aggregate (the
verifier's share is recoverable, the rework share is not), so we do not report a seconds-based
recomputation of $\D$; the conversion band in Section~\ref{sec:results} carries the unit caveat
instead.

Because $\D$ has $\SV$ in the denominator, differences in instance size across families would
produce differences in $\D$ that have nothing to do with verification. Instance sizes were
therefore matched by construction: \emph{per-product} $\SV$ ranges from \VDsmin{} to \VDsmax{}
tokens across the four families, a ratio of \VDsratio{}. (Table~\ref{tab:main} reports $\SV$ summed
over a family's baseline products; the per-product figure is that sum divided by the cell size.)
We report this matching rather than leaving it implicit, because unmatched denominators are the
most likely way for a study of this kind to manufacture its own result.

\subsection{Pre-registration}
\label{sec:prereg}

The success bars, the verifiability rubric and its scores, the $\SV$ baseline procedure, the
decoding settings, the rework cap and the severity rules were written into the repository's
\texttt{README.md} and committed before any debt ratio was computed. The bars were:

\begin{itemize}
\item \textbf{B1.} At least one family shows $\D\geq 1.0$ under the baseline configuration.
\item \textbf{B2.} $\D$ orders monotonically with the independently rated verifiability score
across the four families (Spearman $\leq-0.8$, exact permutation $p\leq0.10$).
\item \textbf{B3.} Guardrails cut $\D$ by at least 30\% on the worst baseline family at no more
than 15\% throughput cost.
\end{itemize}

A missed bar is reported as a finding with its measured value. Rubrics, caps, baselines, severity
rules and the bars themselves (B1--B3) were fixed before the run and were not changed after seeing
results. One change \emph{was} made in response to a result, and we label it as such rather than
fold it into the pre-registration: after the first sweep (v1, 340-token producer budget) returned a
zero first-pass acceptance, we raised the producer budget to 1536 and re-ran (v2, the run reported
throughout). This is a post-hoc protocol amendment. The v1 result --- zero first-pass census
acceptance (0 of 280 baseline) and $\D=5.53$--$8.87$ (worst family invoices, 8.87) --- is reported
as headline metrics only in Table~\ref{tab:versions}. The v1 raw logs were overwritten when the
truncation amendment re-ran the protocol from scratch (the old \texttt{results/raw/*.jsonl} were
archived to \texttt{results/raw\_trunc340\_archive/}), so v1 per-family $\VV/\SV$ and $\RV/\SV$
cannot be recomputed or verified and are not tabulated. The deviation is recorded in
\texttt{README.md}\,\S1.7, but only the v2 logs are retained in the repository. Other post-hoc
analyses are labelled exploratory.

\section{The verification-effort proxy}
\label{sec:proxy}

This section exists because the central measurement in this paper is a proxy, and a paper that
does not say so is asking to be misread.

\paragraph{What is measured.} $\VV$ is the number of tokens an independent verifier agent processes
while checking a product to the rubric. $\RV$ is the number of tokens repair calls process. Neither
is human verification effort. A human checking six invoices opens a spreadsheet, joins on the
purchase order, and looks at a column of differences; the token count of a language model
re-deriving the same comparison from prose is a different quantity with a different cost structure.

\paragraph{Two biases, in opposite directions.} We expected agent token volume to \emph{understate}
human verification effort, and to understate it most where verifiability is lowest: a human with a
tool is cheap per mechanical check, an agent re-reads the whole instance on every pass regardless,
so its token count rises less than a human's time as verifiability falls. That bias would make
measured $\D$ a lower bound on a human $\D$. A second bias points the other way: when the verifier
is wrong it triggers rework a correct checker would not, so $\RV$ prices some of the checker's
error. The positive control bounds this (it rejects \VDgoldRej{} of census-perfect gold instances,
\VDgoldRejRate{}; 95\% CI \VDgoldRejLo--\VDgoldRejHi), and Section~\ref{sec:results} quantifies it
directly: we attribute to the checker only the \emph{excess} over a one-pass gold baseline, because
the first verification pass is a structural floor (below), not error. For every trajectory whose
first rejection was false, excess $= (\VV+\RV) - \SV\cdot\D_{\text{gold}}$ (matched to the
trajectory's family and model, and floored at zero, since no trajectory costs less than one
verification pass), summed over those trajectories and bootstrapped over products ($B=4000$, seed
13; the gold baseline is a fixed positive control and is \emph{not} resampled). The checker-attributable
share is \VDchkShare{} of all $\VV+\RV$ (95\% CI \VDchkShareLo--\VDchkShareHi; excess
$\D_{\text{check}}=\VDchkD$), so the checker-error component inflates $\D$ but does not dominate it.
Measured $\D$ is therefore a measurement of this pipeline, with the checker-error increment measured
rather than assumed.

\paragraph{The structural floor from the accounting.} $\VV$ is the number of tokens the verifier
\emph{processes}: system prompt, rubric, raw inputs, the product under review, and the verdict. $\SV$
prices only the raw inputs and the canonical deliverable, with no prompt scaffolding. A single
verification pass therefore already costs about what $\SV$ costs, before any rework. The gold
control makes this concrete: a perfect product, verified once, with zero rework, has
$\D_{\text{gold}}=\VDgoldDmin$--$\VDgoldDmax$. \textbf{B1 is therefore guaranteed by the unit asymmetry,
not discovered by the experiment}: any pipeline whose verifier must re-read the inputs will show
$\D\gtrsim 1$ even when every product is correct. The floor is a property of this protocol
--- its rubric, input representations and verifier configuration --- not a universal constant
of verification. We report B1 as met but say plainly that it is a
structural check, not evidence about agent work; the informative quantities are the excess over
$\D_{\text{gold}}$ and the verifier's error rate.

\paragraph{What we do instead of converting.} We never convert agent time into human hours or
dollars. Instead we report the whole family of ratios implied by a band of conversion factors
$\lambda$, where $\lambda$ multiplies measured agent effort before it is compared with $\SV$:
\begin{equation}
\D(\lambda) \;=\; \frac{\lambda(\VV+\RV)}{\SV}.
\label{eq:lambda}
\end{equation}
Since $\SV$ does not move, $\D$ is linear in $\lambda$: the \emph{level} of every debt ratio scales
with the conversion factor, while the \emph{ordering} across families is invariant. The break-even
is $\lambda^{*}=1/\D$: the ratio stays above one for any conversion factor above
$\lambda^{*}=\VDlambdaStarMin$--$\VDlambdaStarMax$ across families. That is a more useful statement
than the band alone --- it says how far the token-to-human scaling could be wrong before the
qualitative conclusion changes --- and it is the form in which the result should be read. Section
\ref{sec:robust} reports $\D(\lambda)$ for $\lambda\in\{0.5,1,2,4\}$.

\paragraph{Why tokens and not a common human unit.} A ratio needs a common unit, and the obvious
candidate --- human minutes --- is undefined when the production cost approaches zero, which is the
case the construct is about \citep{crescitelli2026}. Tokens are well defined for both the
production saving and the checking effort, so we use them and state the conversion explicitly
rather than hiding it in a minutes-per-token constant. The cost is that the level is a proxy; the
benefit is that the quantity is measurable at all.

\section{Experimental setup}
\label{sec:setup}

\paragraph{Models and decoding.} Two open-weight models served locally by Ollama on a single
workstation act as \emph{complete agents} that both produce and self-verify their work:
\texttt{qwen2.5:14b-instruct} (primary) and \texttt{llama3.1:8b} (secondary). Decoding is fixed per
role and logged with every artifact: temperature $0$, top-$p$ $1$, fixed seed, context window
4096, a per-role generation budget of 1536 tokens for the producer and rework and 900 for the verifier, and a
mild repetition penalty of 1.05 (the 8B model occasionally loops while emitting its verdict JSON;
appendix~\ref{app:repeats} documents the handling). No API keys and no hosted endpoints are used
anywhere in the code; the backend refuses any non-loopback host.

\paragraph{Scale.} \VDnBase{} work products per family per model in the baseline arm, and
\VDnGuard{} per family per model in the guardrail arm, for a total of \VDnumRecords{} records. Each
record stores the seed, the model, the decoding configuration, the produced text, the verifier's
verdict, the census result at first pass and at the end, and the token and wall-clock accounting
for production, verification and rework separately.

\paragraph{Analysis.} The primary estimate of $\D_f$ is the ratio of sums over the family's
instances --- total cost per unit of total work displaced --- with a 4000-draw instance bootstrap
for the interval; the mean of per-product ratios is reported alongside. The rank correlation uses
the pre-registered verifiability score with an exact permutation $p$-value over all $4!=24$
orderings. Guardrail deltas are computed on the same instances in both arms.

\section{Results}
\label{sec:results}

\subsection{Debt ratios}

Table~\ref{tab:main} reports the decomposition. Every family carries a debt ratio far above one:
on all four families and both models, the cost of checking and repairing what the agent produced is
a multiple --- 5.7 to 11.7, depending on family and model --- of the production saving. \textbf{B1
is \VDbarB} by its letter, but the positive control below shows that this is a property of the
accounting, not a discovery, and we report it that way. The worst family is \VDworstFamily{}, at
$\D=\VDWorstD$ (95\% interval $\VDWorstDlo$--$\VDWorstDhi$).

\paragraph{What the checker catches, and what it misses.} With the 1536-token producer budget the
producer does ship work the census accepts: on \famname{cleansing}, \VDCleansePass{} of products are
free of critical and major defects at first pass and \VDCleanseFinal{} after rework, while the other
three families remain at zero (Table~\ref{tab:main}). Matching the verifier's first-pass decision to
the census at the product level --- the level at which accept/reject is made, with \emph{positive}
defined as the census flagging a blocking defect --- gives \VDdetTP{} true positives,
\VDdetFP{} false rejections, \VDdetFN{} false accepts and \VDdetTN{} true negatives, i.e.\ precision
\VDdetPrec{} and recall \VDdetRec{}. The error is not diffuse; it is concentrated exactly where clean
work exists. On the three families whose producer never clears the bar the verifier's precision is
essentially perfect (invoices \VDdetInvoicesPrec{}, emails \VDdetEmailsPrec{}, reports
\VDdetReportsPrec{}); on \famname{cleansing} \VDdetCleanseFalseRej{} of its rejections are false ---
it accepted \VDCleanseVerAcc{} of the \VDdetFP{} census-clean products (precision
\VDdetCleansePrec{}) --- and on census-perfect gold it rejects \VDgoldRej{}. Attributing only the excess over a one-pass gold baseline (\S\ref{sec:proxy}) of every product
whose first rejection was false --- verification \emph{and} rework beyond the unavoidable first
pass ---
the checker-attributable share is \VDchkShare{} of all $\VV+\RV$ (95\% CI \VDchkShareLo--\VDchkShareHi),
or an excess $\D_{\text{check}}=\VDchkD$ (95\% CI \VDchkDLo--\VDchkDHi) against a total
$\D=\VDchkDTotal$; so checker error inflates $\D$ but is a minority of it, and the bulk of the work
repairs genuine defects. The verifier is a reject-biased detector whose precision collapses as the clean base rate
rises: it is right on dirty work and wrong on clean work, and only \famname{cleansing} and the gold
control contain enough clean work to expose it. On the three families where the producer still ships
defective work at first pass, rework ends at its stall detector for the large majority of items
(87--100\% depending on family), with the remainder hitting the five-round cap; that is why the
measured $\RV$ is the cost of \emph{attempted} repair rather than of successful repair.

\paragraph{The positive control.} To separate the accounting from the work, we ran the planted gold
deliverable --- the correct product, by construction --- through the same verifier and census. The
census accepts it on every instance (0/5 blocking, all families), so the gold is not the problem.
The verifier rejects it anyway: across both models it rejects the planted gold on
\VDgoldRej{} of instances (\VDgoldRejRate{}; 95\% CI \VDgoldRejLo--\VDgoldRejHi; per-family and
per-model rates in Table~\ref{tab:gold}). With one verification pass and zero rework, the ratio is
already $\D_{\text{gold}}=\VDgoldDmin$--$\VDgoldDmax$. The checker is not uniformly bad, and the
split is the informative part: on the \emph{dirty} side it performs well, rejecting \VDverCaught{}
of the \VDverBlocked{} census-blocked baseline products (recall \VDverRecall); it is the
\emph{clean} side where it fails, rejecting \VDgoldRej{} of the gold instances (\VDgoldRejRate{};
95\% CI \VDgoldRejLo--\VDgoldRejHi). That the
verifier's acceptance of gold (\VDgoldAcc{}) is statistically indistinguishable from its acceptance
of production work (\VDBaseAccPooled{} pooled over both models; Fisher exact $p=\VDgoldFisher$, not
significant) means the checker shows essentially no discrimination between clean and defective
products. We
therefore \emph{cannot} read the production rework as dominated by false positives: the census
flagged all but \VDfirstPass{} of the baseline products at first pass, so most rejections are correct
catches (precision \VDdetPrec{}). The control and the product-level agreement above do more than
assert over-rejection --- they quantify it: \VDdetFalseRej{} of first-pass rejections are false,
concentrated on \famname{cleansing} (\VDdetCleanseFalseRej{}) and on gold (\VDgoldRej{}), and
the excess over a one-pass gold baseline of trajectories with a false first rejection is
\VDchkShare{} of all $\VV+\RV$ (95\% CI \VDchkShareLo--\VDchkShareHi). Two consequences follow, and
they are the honest reading of the headline numbers: (i) B1 is guaranteed by the unit asymmetry in
Section~\ref{sec:proxy}, since a verifier that re-reads the inputs costs about what $\SV$ prices;
and (ii) the excess of the measured $\D$ over $\D_{\text{gold}}$ is now decomposed into a measured
checker-error component (\VDchkD{} of $\D$) and genuine defect repair, not assumed
wholesale; a stronger or human checker would move both $\VV$ and $\RV$, and the false-reject rate,
not the raw level, is the quantity to watch.

\begin{table}[t]
\centering
\small
\caption{Positive control: the planted gold deliverable through verifier and census, both models,
5 instances per family. The census accepts every gold product; the verifier mostly does not.
False-reject is the share of gold the verifier rejects; pooled false-reject \VDgoldRej{}
(\VDgoldRejLo--\VDgoldRejHi, 95\% CI).}
\label{tab:gold}
\begin{tabular}{llrrr}
\toprule
family & model & gold accepted & false-reject & $\D_{\text{gold}}$ \\
\midrule
\famname{invoices}  & Qwen  & \VDGoldAccqwenInvoices & \VDGoldFrqwenInvoices & \VDGoldDqwenInvoices \\
\famname{invoices}  & Llama & \VDGoldAccllamaInvoices & \VDGoldFrllamaInvoices & \VDGoldDllamaInvoices \\
\famname{cleansing} & Qwen  & \VDGoldAccqwenCleanse & \VDGoldFrqwenCleanse & \VDGoldDqwenCleanse \\
\famname{cleansing} & Llama & \VDGoldAccllamaCleanse & \VDGoldFrllamaCleanse & \VDGoldDllamaCleanse \\
\famname{emails}    & Qwen  & \VDGoldAccqwenEmails & \VDGoldFrqwenEmails & \VDGoldDqwenEmails \\
\famname{emails}    & Llama & \VDGoldAccllamaEmails & \VDGoldFrllamaEmails & \VDGoldDllamaEmails \\
\famname{reports}   & Qwen  & \VDGoldAccqwenReports & \VDGoldFrqwenReports & \VDGoldDqwenReports \\
\famname{reports}   & Llama & \VDGoldAccllamaReports & \VDGoldFrllamaReports & \VDGoldDllamaReports \\
\bottomrule
\end{tabular}
\end{table}

\begin{table}[t]
\centering
\small
\caption{Debt-ratio decomposition, pooled over both models. $\SV$, $\VV$, $\RV$ in thousands of
tokens; $\D=(\VV+\RV)/\SV$; interval is a 4000-draw instance bootstrap. ``pass'' is the share of
products with no critical or major defect at first pass; ``final'' after rework.}
\label{tab:main}
\begin{tabular}{lrrrrrrrr}
\toprule
family & $\SV$ & $\VV$ & $\RV$ & $\VV/\SV$ & $\RV/\SV$ & $\D$ & pass & final \\
\midrule
\famname{invoices}  & \VDInvoicesS & \VDInvoicesV & \VDInvoicesR & \VDInvoicesVS & \VDInvoicesRS & \VDInvoicesD{} [\VDInvoicesDlo, \VDInvoicesDhi] & \VDInvoicesPass & \VDInvoicesFinal \\
\famname{cleansing} & \VDCleanseS & \VDCleanseV & \VDCleanseR & \VDCleanseVS & \VDCleanseRS & \VDCleanseD{} [\VDCleanseDlo, \VDCleanseDhi] & \VDCleansePass & \VDCleanseFinal \\
\famname{emails}    & \VDEmailsS & \VDEmailsV & \VDEmailsR & \VDEmailsVS & \VDEmailsRS & \VDEmailsD{} [\VDEmailsDlo, \VDEmailsDhi] & \VDEmailsPass & \VDEmailsFinal \\
\famname{reports}   & \VDReportsS & \VDReportsV & \VDReportsR & \VDReportsVS & \VDReportsRS & \VDReportsD{} [\VDReportsDlo, \VDReportsDhi] & \VDReportsPass & \VDReportsFinal \\
\bottomrule
\end{tabular}
\end{table}

Figure~\ref{fig:decomp} shows the same decomposition graphically, with the unit line at $\D=1$.
The visual point is the one the identity makes: the bars are not small relative to the unit. For
\famname{reports}, verification alone costs \VDReportsVS{} times the production saving, before any
rework is counted.

\begin{figure}[t]
\centering
\includegraphics[width=0.52\linewidth]{figures/fig1_decomposition.pdf}
\caption{Debt-ratio decomposition by family, primary model. Verification (dark) and rework
(light) cost per unit of production saving. The dashed line is $\D=1$, where the deployment is
value-neutral.}
\label{fig:decomp}
\end{figure}

\subsection{The verifiability gradient}

Figure~\ref{fig:monotone} plots $\D$ against the verifiability score assigned before any run.
\textbf{B2 is \VDbarC}, and it fails by showing no detectable gradient rather than a clean
inversion. The pre-registered hypothesis was that debt rises as verifiability falls. The measured
Spearman correlation is $\rho_s=\VDspearman$ with exact permutation $p=\VDpperm$: the sign is
positive, but with four families and $p=0.75$ the data do not support a claim of reversal --- they
support the weaker statement that the pre-registered gradient is not there. The correlation is flat
on the secondary model ($\rho_s=\VDqSpearman$ primary, $\rho_s=\VDlSpearman$ secondary), so it is
not a pooled artifact of one model. Note also that with $n=4$ the pre-registered threshold
($\rho_s\le-0.8$, $p\le0.10$) is effectively a requirement of perfect monotonicity: only
$\rho_s=-1$ clears $p\le0.10$. B2 was a weak test by construction, and we say so rather than
over-read its failure.

We can say something about \emph{why} the gradient is absent, but it is a conjecture, and the
obvious version of it is contradicted by our own table. Where work is verifiable, the checker can
localise a defect to a field, so the loop has something to act on: it re-derives, repairs,
re-verifies, and books the tokens. But the number of graded checks does not by itself predict
$\D$: \famname{cleansing} has the most checks per product (\VDchecksCleanse) and a lower ratio
($\VDCleanseD$) than \famname{invoices} (\VDchecksInvoices{} checks, $\VDInvoicesD$) or
\famname{reports} (\VDchecksReports{} checks, $\VDReportsD$). What seems to separate the families is
whether the checker can point to a \emph{specific wrong field}, not how many fields exist; that is
consistent with the positive control (the verifier rejects correct work, so its false positives,
not the check count, drive rework), but we have not measured localisability directly. The
defensible claim is narrower than the one we would like: the ledger records what the checking
process can see and act on, and on this data that quantity is not ordered by verifiability.

\begin{figure}[t]
\centering
\includegraphics[width=0.48\linewidth]{figures/fig2_monotone.pdf}
\caption{Debt ratio against pre-registered verifiability score, pooled over models, with bootstrap
intervals. The dotted line is a least-squares fit shown for orientation only; the test is the rank
correlation.}
\label{fig:monotone}
\end{figure}

\subsection{Guardrails as debt management}

The guardrail arm adds two things and nothing else: a structured acceptance checklist the producer
must complete before emitting its product, and a mandated second verification pass before delivery.
Figure~\ref{fig:guard} and Table~\ref{tab:guard} report the trade.

\begin{table}[t]
\centering
\small
\caption{Guardrail arm, compared on the same instances in both arms (matched $\VDnGuard$ per
family per model). $\Delta\D$ is the change in the debt ratio; throughput cost is the percentage
increase in total effort (production plus verification plus rework) per work product.}
\label{tab:guard}
\begin{tabular}{lrrrrr}
\toprule
family & $\D$ baseline & $\D$ guarded & $\Delta\D$ (\%) & throughput cost & pass (guarded) \\
\midrule
\famname{invoices}  & \VDInvoicesGuardBase & \VDInvoicesGuardD & \VDInvoicesGuardChange & \VDInvoicesGuardCost & \VDInvoicesGuardPass \\
\famname{cleansing} & \VDCleanseGuardBase & \VDCleanseGuardD & \VDCleanseGuardChange & \VDCleanseGuardCost & \VDCleanseGuardPass \\
\famname{emails}    & \VDEmailsGuardBase & \VDEmailsGuardD & \VDEmailsGuardChange & \VDEmailsGuardCost & \VDEmailsGuardPass \\
\famname{reports}   & \VDReportsGuardBase & \VDReportsGuardD & \VDReportsGuardChange & \VDReportsGuardCost & \VDReportsGuardPass \\
\bottomrule
\end{tabular}
\end{table}

\begin{figure}[t]
\centering
\includegraphics[width=0.48\linewidth]{figures/fig3_guardrail.pdf}
\caption{Guardrail arm: increase in $\D$ (positive = debt rose) against throughput cost. The
dashed lines mark the pre-registered B3 corner (a 30\% cut at $\le$15\% cost).}
\label{fig:guard}
\end{figure}

\textbf{B3 is \VDbarD} --- and it fails in the direction opposite to the intervention's intent.
Compared on the \emph{same} instances in both arms, the guardrail did not cut $\D$; it raised it,
from $\D=\VDWorstGuardBase$ to $\D=\VDWorstGuardD$ on the worst family, a change of
$+\VDWorstGuardDlo$ to $+\VDWorstGuardDhi$ (95\% CI, paired bootstrap) at a throughput cost of
\VDWorstGuardCost. The pattern holds on every family (Table~\ref{tab:guard}): the debt ratio rises
by \VDguardRiseLo--\VDguardRiseHi and total effort per product rises by
\VDguardCostLo--\VDguardCostHi, with paired intervals that exclude zero on all four families.

The decomposition says where the money went. The mandated second verification pass roughly doubles
$\VV$ by construction --- the second pass is billed in full --- and the checklist was supposed to
buy that cost back by preventing rework. Mostly it did not: first-pass blocking defects barely move
(from about 3.9 to 3.7 per product on \famname{invoices}), the verifier still accepts essentially
nothing (see the accept columns in Table~\ref{tab:main}), and mean rework rounds actually \emph{rise}
on three of four families. There is one honest counterexample: on \famname{cleansing} the guardrail
did raise the census pass rate, from \VDCleansePass{} to \VDCleanseGuardPass{} (about nine
percentage points), so the checklist is not inert. But the gain does not pay for itself: $\D$ on
\famname{cleansing} still rises from \VDCleanseGuardBase{} to \VDCleanseGuardD{}, because the second
pass is billed in full and the acceptance gain is too small to offset it. On the three families where
the producer never clears the bar the guardrail is a pure cost adder, and the ratio registers it as
such.

This failure is partly a design consequence, and we say so. With the verifier accepting
essentially nothing in both arms, the guardrail's only acceptance gain is the \famname{cleansing}
pass-rate improvement, which is too small to offset the doubled verification cost; and because the
second pass doubles $\VV$ while $\RV$ is truncated by the stall detector, $\D$ can only rise.
The result is still informative --- it shows that a verification mandate applied to a producer that
mostly never clears the bar is mostly a cost, with a small acceptance gain only where the producer
already ships clean work --- but it is a weaker test of the instrument than the pre-registration
implied, and a stronger producer is needed to give the guardrail a fair chance
(Section~\ref{sec:limitations}).

\subsection{Leakage: what the checker does not see}
\label{sec:leakage}

The rework loop is gated by the verifier, not by ground truth, because that is what a deployment
does. The gap between the two is where defects ship. On this experiment the gap is degenerate in
an instructive way: the verifier accepted only \VDqBaseAcc{} products on the primary model and
\VDlBaseAcc{} on the secondary, and the primary accepts that occurred still carried a critical
defect under the census (\VDqLeak{} leak rate). That 100\% is an anecdote, not an estimate: with two accepted products the
interval is uninformative, and we report the count rather than a rate. The informative fact is the
denominator --- a checker that accepts almost nothing produces almost no \emph{accepted} work to
leak, because everything is escalated instead.

Two readings, both uncomfortable. The benign one: almost nothing ships through the accept path,
so leakage of \emph{accepted} work is small here --- the failure mode is stalled escalation, not
silent delivery. The uncomfortable one: a checking process this strict is one parameter change
away from a lax one. Relax the verifier's rubric slightly and its acceptance rate rises; the
products it waves through are exactly the ones the census shows are defective, and the leak rate
on this data was 100\% of what it did accept. $\D$ measured under a weak checker is therefore
not comparable to $\D$ measured under a strong one: the weak checker books less $\VV$, triggers
less $\RV$, and delivers the defects. That is why the protocol measures the deterministic census
alongside the loop, and why $\RV$ should be read as rework under the workflow's own checking
process, not under an ideal one. Nor can the bill be cut by staffing the check with more copies
of the producing model: a companion study measures difficulty-adjusted failure correlation of
$\phi = 0.80$ between same-model agents, so a same-model checker inherits the producer's
failures rather than policing them \citep{wang2026onebrain}.

\subsection{Robustness}
\label{sec:robust}

We check four things: the extended ratio that includes production cost, the stability of the
ordering across models, the unit, and the conversion band. All four are reported in
Appendix~\ref{app:robust}. The short version: the extended ratio raises every level and leaves the
pooled ordering unchanged; the ordering is \emph{not} stable across the two models, so only the
level replicates; and the break-even conversion is
$\lambda^{*}=\VDlambdaStarMin$--$\VDlambdaStarMax$ (pooled over both models), which is the most
robust statement in the paper because it survives any plausible human-to-agent scaling.

\paragraph{Loop termination is a measurement choice, not a finding.} The rework loop stops at a
cap of five rounds or when its stall detector trips (the same flagged field set twice, or two
consecutive rounds with no reduction in the defect count), so the measured $\RV$ is the cost of
\emph{attempted} repair under that rule --- a reviewer rightly asked whether the debt levels are an
artifact of the rule. We replay the stored per-round verdicts under tighter rules to bound the
dependence: on the worst family (pooled over both models) the debt ratio is $\D=\VDtermCapA$ at a
one-round cap, rising to $\D=\VDtermCapE$ at the deployed five-round cap --- the same
$\VDWorstD$ reported in Table~\ref{tab:main}, which the replay reproduces exactly (a unit test
asserts this) --- and the stall detector (patience~2) stops most products before the cap, so the cap
mainly bounds the worst case. The level therefore scales with the rule; what does \emph{not} scale is
the sign (debt above one) and the structural floor, which hold under every rule we replayed.

\section{Discussion}
\label{sec:discussion}

\paragraph{The accounting implication.} The identity in Equation~\eqref{eq:identity} is not a
critique of agents; it is a statement about what a partial sum does. A firm that books $\SV$ and
not $\VV+\RV$ is not lying, and its production numbers are not wrong. It is reporting a gross
figure as if it were net. The remedy is not to stop deploying. It is to book the second line, and
the contribution of this paper is to make the second line a measurable quantity with a name and a
unit rather than an anecdote about review queues.

\paragraph{Reading the field pattern through the construct.} The coding-agent evidence
\citep{chen2026} is consistent with a debt ratio near one in a high-verifiability family: output up
30\%, review time up 49\%, delivered work flat. Our construct says what that pattern means and
what would have to be true for it to hold: verification cost rose roughly in proportion to the
production saving, and rework absorbed the remainder. Our own high-verifiability family
(\famname{invoices}) sits far above that, at $\D\approx\VDqDInvoices$, but the positive control shows why:
our verifier rejects most correct work and the accounting has a structural floor, so the lab number
is not a field estimate. The study cannot tell us how the pattern moves when reviewers are people
rather than a same-model verifier; that is the field measurement this paper does not attempt.

\paragraph{Guardrails are debt service, not debt avoidance.} The intervention result is best read
as a portfolio choice, and on these configurations the choice is priced and rejected. A checklist
plus a second pass raises $\VV$ by construction; whether it lowers $\D$ depends on whether the
reduction in $\RV$ exceeds the increase in $\VV$. Here it did not: the checklist left the
first-pass defect rate essentially intact, the second pass doubled verification, and $\D$ rose on
every family. That is not a failure of the construct but the construct doing its job --- pricing a
quality intervention instead of assuming it. The framing matters: a guardrail is not free quality,
it is a payment, and the right question is the price. A deployment that cannot reach its
acceptance standard should be fixed at the source (a stronger producer, or a smaller and more
checkable task), not asked to pay for more checking of output that will be rejected anyway.

\paragraph{The verifiability result, read correctly.} The pre-registered ordering did not appear,
and the honest reading is narrower than the one we pre-registered: with four families and
$p=0.75$, the data show no detectable gradient, not a reversal. What the experiment does support is
a claim about the \emph{object} being measured. The ratio prices cost booked by a checking process,
and the positive control shows that process rejecting correct work most of the time, so part of
the measured ratio is the checker's error --- now quantified, as the excess over a one-pass gold
baseline, at \VDchkShare{} of all $\VV+\RV$ (95\% CI \VDchkShareLo--\VDchkShareHi) --- while the
census shows the baseline
work was mostly genuinely defective, so most rejections are correct catches (precision \VDdetPrec{}).
Two deployments can therefore report very
different $\D$ for work of the same true quality, and the one reporting the \emph{lower} ratio may
be the one shipping the worse product. This is why the construct is a measurement protocol, not a
single number: the ratio without the census and the control beside it is a statement about the
checker, not about the work. It also predicts a testable field pattern: organisations with weak
review will report attractive agent ROI and high defect escape rates at the same time.

\paragraph{Why the ratio and not the defect rate.} A defect rate tells you how often the agent is
wrong. It does not tell you what finding out costs, and it does not tell you what repairing costs.
Two families with identical defect rates can have very different debt ratios if one's defects are
cheap to detect and the other's are not. The ratio is the quantity that enters the ROI
calculation; the defect rate is an input to it.

\section{Limitations}
\label{sec:limitations}

\paragraph{No field measurement.} This paper measures no organisation's ROI. Every number here
comes from a laboratory with synthetic tasks and planted truth. The ROI implications are stated as
accounting identities plus measured lab ratios --- ``if a family's field $\D$ resembles the lab
$\D$, net gain is $\SV(1-\D)$, which is negative for every ratio we measured'' --- and nothing
stronger. A field study would need to instrument verification and rework in comparable units,
which is possible but is not this paper.

\paragraph{A single rater for verifiability.} The verifiability scores were assigned from the
rubric before any run, but by a single rater, so no inter-rater agreement is reported; the README
sketches a second-rater protocol that this run did not execute. The scores should be read as one
analyst's prior, and the B2 failure is correspondingly less about the scores than about the absence
of any gradient to score.

\paragraph{The verifier's error is concentrated on clean work.} The positive control
(Section~\ref{sec:results}) is not a detail: the same-model verifier rejects \VDgoldRej{} of the
planted gold instances (\VDgoldRejRate{}; 95\% CI \VDgoldRejLo--\VDgoldRejHi). That demonstrates
over-rejection of clean work, and the product-level agreement above quantifies it: \VDdetFalseRej{}
of first-pass rejections are false, concentrated on \famname{cleansing} (\VDdetCleanseFalseRej{})
and on gold, and the excess over a one-pass gold baseline of checker-attributable trajectories is
\VDchkShare{} of all $\VV+\RV$ (95\% CI \VDchkShareLo--\VDchkShareHi). It does \emph{not}
establish that production rework is dominated by false positives: the census flagged all but
\VDfirstPass{} of the baseline products, so most production rejections are correct catches
(precision \VDdetPrec{}), and the verifier's precision collapses only where clean work is common. A
same-model verifier is also the worst case
for correlated error \citep{huang2023}: it shares the producer's failure modes, so its rejections
and the producer's defects are not independent. A stronger or human verifier would lower
both $\VV$ and $\RV$ and change every level in this paper; the direction of the qualitative
conclusions (debt above the structural floor, a guardrail that costs without curing) is what we
would expect to survive, but we have not measured it.

\paragraph{Capacity.} For this revision the producer and rework budget is 1536 tokens and the
verifier budget 900, with context 4096. A reviewer concern was that the original 340-token producer
budget truncated many products mid-JSON and could have manufactured a zero first-pass acceptance;
we raised the budget and re-ran --- a post-hoc amendment, labelled as such in \S\ref{sec:prereg}.
The truncation concern was real for the producer: with room to finish, it now ships census-clean
work at first pass on \famname{cleansing} (\VDCleansePass{}), so the zero first-pass census
acceptance in v1 was in part a truncation artefact. What stays near zero is the \emph{verifier's}
acceptance (\VDCleanseVerAcc{} on \famname{cleansing}; \VDqBaseAcc{} primary, \VDlBaseAcc{}
secondary). But the token cost of that over-rejection is bounded: the excess over a one-pass gold
baseline of checker-attributable trajectories is \VDchkShare{} of all $\VV+\RV$ (95\% CI
\VDchkShareLo--\VDchkShareHi), so the acceptance
bottleneck is the verifier's, while the debt level is mostly the cost of repairing genuine defects.

\paragraph{External validity of synthetic tasks.} The four families are recognisable classes of
process work, but they are generated, and generated work is cleaner than real work: the inputs are
well formed, the rules are stated, and the ground truth is exact. Real process work has ambiguous
inputs, undocumented exceptions, and standards that live in people's heads. All of those raise
verification cost, so the lab ratios are more likely to be too low than too high --- but that is
an argument, not a measurement.

\paragraph{An acceptance standard the \emph{verifier} never meets.} The census acceptance rule is
deliberately strict --- zero critical or major defects --- and on three of four families neither
model's producer met it on a single product; on \famname{cleansing}, by contrast, \VDCleansePass{}
met it at first pass and \VDCleanseFinal{} after rework. What never clears the bar is the
\emph{verifier}: it accepts \VDqBaseAcc{} of baseline products on the primary model and
\VDlBaseAcc{} on the secondary, and it rejects census-clean \famname{cleansing} output
(\VDCleanseVerAcc{} accepted) along with \VDgoldRej{} of census-perfect gold. That makes the level of
$\D$ a statement about a pipeline whose verifier always escalates, not about one that sometimes
succeeds, and it removes the acceptance-rate channel from the guardrail comparison (the verifier
accepts essentially nothing in either arm, so the guardrail can only be judged on cost). A looser
standard, or a stronger \emph{verifier}, would move both $\VV$ and $\RV$ down and would let the
checklist show whatever benefit it has. The bars and the rubric were fixed before the run and are
not loosened here; the limitation is stated so the numbers are read at the right altitude.

\paragraph{Judge and rubric dependence.} The census is deterministic and therefore reproducible,
but it is only as good as the schema it grades: a product can be substantively right and
schematically wrong, and the census will call it wrong. The verifier is the weaker link, and the
positive control prices it (Section~\ref{sec:results}); a stronger checker would change both
$\VV$ and $\RV$.

\paragraph{One workstation, two models.} Both models are 8--14B open-weight models served on one
machine. Larger models would produce better first drafts and probably cheaper verification;
smaller ones the reverse. $\D$ is indexed by configuration for exactly this reason, and the numbers
here should be read as a measurement of these configurations, not of agents in general.

\paragraph{The proxy.} Section~\ref{sec:proxy} states the direction of the bias and the
identification limits. The short version: levels are not identified, orderings are, and the
ordering is tested on a compressed version of the gradient.

\section{Conclusion}

The cost of agent-produced work does not disappear when the production step gets cheap. It moves
to the step that decides whether the work is any good, and to the step that fixes it when it is
not. This paper gives that movement a unit of account: three costs, one ratio, one identity. The
most important result is not the level of the ratio but what the positive control did to it. The
ratio is above one even for a perfect product, because a verifier that re-reads the inputs costs
about what the saving prices; and it is far above that floor partly because the checker rejects
correct work, as the positive control demonstrates. So the honest summary is: the accounting
has a structural floor, part of the measured excess is checker error --- \VDdetFalseRej{} of
first-pass rejections are false, and such trajectories are \VDchkShare{} of all $\VV+\RV$ --- and
the protocol is what lets a reader see both.

The pre-registered results failed, and the failures are the useful part. The verifiability
ordering did not appear, and we now read the ratio as a property of the checking process rather
than of the work; the guardrail raised debt, because it billed a second verification pass without
preventing the rework it was meant to prevent. Together with the positive control, these say
something specific and testable: a deployment's measured debt is dominated by how reliable its
checker is, so two organisations with the same true quality can report very different ratios, and
the one reporting the lower number may be shipping the worse product.

The practical instruction is short. When an agent deployment is proposed, ask for three numbers,
not one: the production saving, the cost of finding out whether the output is right, and how often
the checker is wrong on work that is known to be correct. If any of the three is missing, the
business case is a partial sum. The construct and the control in this paper are what make the
second and third numbers measurable.

\appendix

\section{Robustness details}
\label{app:robust}

\paragraph{The extended ratio.} Including production cost, $\D^{\text{full}}=(P+\VV+\RV)/\SV$,
raises every ratio (\VDdfullMin--\VDdfullMax pooled) and leaves the pooled family ordering unchanged. The
reported $\D$ is therefore conservative, as Section~\ref{sec:construct} said it would be.

\paragraph{Model heterogeneity.} The pooled ordering is not stable across models and should not
be over-read. On the primary model the ordering is \famname{invoices} (\VDqDInvoices) $>$
\famname{reports} (\VDqDReports) $>$ \famname{emails} (\VDqDEmails) $>$ \famname{cleansing} (\VDqDCleanse), with the two
weakest ratios nearly tied; on the secondary model it is \famname{cleansing} (\VDlDCleanse) $>$
\famname{reports} (\VDlDReports) $>$ \famname{invoices} (\VDlDInvoices) $>$ \famname{emails} (\VDlDEmails), again with the two
weakest nearly tied. What replicates across both models is the level --- every family above one by a
multiple --- not the ranking. Rank statements about verifiability need more models than this
experiment has.

\paragraph{Units.} The ratio is expressed in tokens. Wall-clock was logged at the product level
and for each verifier pass, but not separately for rework calls, so a complete seconds-based
recomputation of $\D$ is not possible from the released logs; we do not claim one. The unit caveat
is instead handled by the conversion band, under which the level moves and the ordering cannot
move, because $\D(\lambda)$ is linear in $\lambda$ and $\SV$ does not scale with $\lambda$. The
asymmetry between the two sides of the ratio is itself measurable: on the baseline run,
\VDpromptShareMin--\VDpromptShareMax{} of the tokens the checker processes are \emph{prompt} tokens
--- re-reading the inputs, the rubric and the product --- rather than generated verdicts. The rubric
and prompt scaffolding have no analogue in $\SV$, which is precisely why $\D_{\text{gold}}$ does not
approach zero.

\paragraph{The proxy band and the break-even.} Figure~\ref{fig:sens} reports $\D(\lambda)$ for
$\lambda\in\{0.5,1,2,4\}$. For the worst family the ratio runs from \VDsensMin{} to \VDsensMax{}.
Because $\D(\lambda)$ is linear and $\SV$ does not move, the ratio crosses one at
$\lambda^{*}=1/\D$, which on the primary model (Figure~\ref{fig:sens}'s caliber) is
$\lambda^{*}=\VDlambdaStarQMin$--$\VDlambdaStarQMax$, and pooled over both models
$\lambda^{*}=\VDlambdaStarMin$--$\VDlambdaStarMax$: the conclusion $\D>1$ survives any
human-to-agent effort conversion above roughly $0.15$, and fails only if human checking is cheaper
than that relative to the agent's token volume. This is
the most robust statement in the paper. It does not, however, repair the verifier-error inflation
documented by the positive control; the break-even addresses the unit, not the checker.

\begin{figure}[t]
\centering
\includegraphics[width=0.46\linewidth]{figures/fig4_sensitivity.pdf}
\caption{Debt ratio under a band of human-to-agent effort conversion factors, primary model. The
ratio is linear in $\lambda$; the ordering across families is invariant.}
\label{fig:sens}
\end{figure}

\section{Task family specifications}
\label{app:families}

Each family is a generator plus a deterministic grader, in \texttt{tasks/}. Table~\ref{tab:spec}
summarises the planted structure.

\begin{table}[h]
\centering
\small
\caption{Family specifications. ``checks'' is the number of graded fields per instance.}
\label{tab:spec}
\begin{tabular}{p{2.1cm}p{3.5cm}p{4.2cm}r}
\toprule
family & instance & planted structure & checks \\
\midrule
\famname{invoices} & 6 three-way-match documents & duplicate id, price drift, quantity drift, tax-rate
drift, missing receipt & 18 \\
\famname{emails} & policy KB (7 rules) + 3 emails & defective item, opened item, closed window, clean
case, chargeback threat; value threshold for escalation & 12 \\
\famname{cleansing} & 12--16 raw records, 5 entities, system of record & duplicate clusters, format
drift, blank field, surname trap, contradiction with system of record & 30 \\
\famname{reports} & 2 tables, 10 weeks $\times$ 6 regions & regional decline, single-week anomaly,
aggregate trend, handling-time drift & 10 \\
\bottomrule
\end{tabular}
\end{table}

\section{Severity rules}
\label{app:severity}
Severity is assigned so that the census prices the error the organisation pays for. In
\famname{invoices}, approving a document that should have been held is critical (money leaves), and
holding a clean document is major. In \famname{emails}, an impermissible action set or a reply that
promises a forbidden outcome is critical; a wrong amount is major; a missing rule citation is
minor. In \famname{cleansing}, a wrong system-of-record field or a lost entity is critical; a
normalisation failure is major. In \famname{reports}, a figure that cannot be reproduced from the
tables is critical; a wrong figure in the findings block is major. Acceptance on the census side
means no critical and no major defect.

\section{Verdict parsing and repetition handling}
\label{app:repeats}

Two decoding artefacts had to be handled before the verifier's verdicts could be trusted, and both
are documented here because either one, left silent, would have inflated the debt ratio.

\paragraph{Token-budget truncation.} The verifier's defect list can exceed a small generation
budget, truncating the JSON mid-object. Before the budget was raised, a truncated verdict parsed to
nothing and \emph{every} product was silently rejected, so the rework loop always ran to its cap and
$R$ --- and therefore $D$ --- was inflated for a purely mechanical reason. The verifier budget is now
900 tokens and the \texttt{accepted} field is emitted first in every family's verdict schema, so the
decision survives truncation even if a defect list does not.

\paragraph{Repetition loops.} The 8B model occasionally enters a repetition loop while emitting its
verdict, repeating one defect object until the budget is exhausted, so the JSON never closes. Two
countermeasures are applied: a mild repetition penalty of 1.05 in the decoding options, and a
salvage parser that recovers the \texttt{accepted} decision and every complete defect object from a
truncated tail. On a probe of 21 instances across all four families and both models, verdict parse
failures fell to zero. A verdict that still cannot be parsed is recorded as a verifier error and
treated as an escalation to a human, not as a rework loop, so it cannot silently inflate $R$.

\section{Reproduction}
\label{app:repro}

\begin{verbatim}
python scripts/selftest.py          # census accepts gold, rejects corruption
python run_all.py sweep --n 35 --n-guard 25 --workers 3
python run_all.py summarize         # -> results/summary.json
python make_figs.py                 # -> figures/*.pdf
python emit_numbers.py              # -> numbers.tex
tectonic paper.tex
\end{verbatim}

The sweep is resumable: records are appended to \texttt{results/raw/*.jsonl} and existing
\texttt{run\_id}s are skipped. Every number in this paper is a macro defined in
\texttt{numbers.tex}, generated from \texttt{results/summary.json}, which is a pure function of
the raw logs; \texttt{CLAIMS.md} maps each claim to its source file.

\begin{sloppypar}
\paragraph{Artifact.} The code, the raw logs (\texttt{results/raw/*.jsonl}), the derived
\texttt{summary.json}, and the regeneration script are released at
\url{https://github.com/boyuw/verification-debt} (a permanent DOI will be
minted). The termination-rule replay is covered by \texttt{tests/test\_termination\_replay.py}, which
asserts that replaying the deployed rule reproduces the reported $\D$ exactly. Mechanically: at the
deployed rule the replay consumes every stored round and sums the stored per-round verifier tokens
and the stored rework \emph{total}, so it is exact; at tighter rules it truncates the stored rounds,
and because the logs record the rework total but not its per-round split, each rework round under a
tighter rule is charged the product's mean rework cost (so tighter-rule levels are approximations,
while the deployed-rule level is not).
\end{sloppypar}

\paragraph{Environment and determinism.} All runs used Ollama~\texttt{0.35.0} (client
\texttt{0.40.6}) on a single Apple M2~Max workstation with 32\,GB unified memory. Models were
\texttt{qwen2.5:14b-instruct} and \texttt{llama3.1:8b}, both the default Ollama pull
(\texttt{Q4\_K\_M} quantization). Decoding was set to greedy: temperature~0, top-$p$~1,
seed~\texttt{13}, context window 4096, producer/rework budget 1536 tokens, verifier budget 900
tokens. Each cell is a single run. We \emph{assume} greedy decoding is bit-reproducible on this
hardware but have \emph{not} verified token-for-token identity across reruns --- batched inference
can differ by hardware and batch composition --- so we report no variance band and flag this as a
reproducibility check still to be run (a 10--20 product rerun would settle it).
The exact model tags, quantization, and Ollama version are reproducible from
\texttt{ollama show <tag>} and \texttt{ollama --version}.

\begin{table}[h]
\centering
\small
\caption{Protocol versions. v1 used a 340-token producer budget; v2 (reported throughout) uses
1536. The change was made after v1 returned zero first-pass acceptance (a post-hoc amendment,
\S\ref{sec:prereg}). v1 per-family $\VV/\SV$ and $\RV/\SV$ were not retained and cannot be
recomputed --- the v1 raw logs were overwritten by the v2 re-run (archived to
\texttt{results/raw\_trunc340\_archive/}) --- so only headline metrics are tabulated (v1 figures as
reported at the time).}
\label{tab:versions}
\begin{tabular}{lrr}
\toprule
metric & v1 (340) & v2 (1536) \\
\midrule
first-pass census pass (baseline)          & 0/280        & 22/280 \\
debt ratio $\D$, range over families       & 5.53--8.87   & 6.62--8.89 \\
worst family (invoices) $\D$               & 8.87         & 8.89 \\
guardrail $\Delta\D$ (rise)                & 43--98\%     & 72--112\% \\
guardrail throughput cost                  & 39--89\%     & 63--101\% \\
verifier accepts (Qwen / Llama)            & not retained & 2/140, 0/140 \\
\bottomrule
\end{tabular}
\end{table}

\begin{thebibliography}{9}
\bibitem[Brynjolfsson(1993)]{brynjolfsson1993}
E.~Brynjolfsson. The productivity paradox of information technology.
\emph{Communications of the ACM}, 36(12):66--77, 1993.

\bibitem[Brynjolfsson and Hitt(1996)]{brynjolfsson1996}
E.~Brynjolfsson and L.~Hitt. Paradox lost? Firm-level evidence on the returns to information
systems spending. \emph{Management Science}, 42(4):541--558, 1996.

\bibitem[Brynjolfsson et al.(2021)]{brynjolfsson2021}
E.~Brynjolfsson, D.~Rock, and C.~Syverson. The productivity J-curve: How intangibles complement
general purpose technologies. \emph{American Economic Journal: Macroeconomics}, 13(1):333--372,
2021.

\bibitem[Chen and Stratton(2026)]{chen2026}
F.~Chen and J.~Stratton. Artificial intelligence in the firm: Bottlenecks in software production.
Working paper (job market paper), Harvard University, August 2026 version.

\bibitem[Cui et al.(2024)]{cui2024}
K.~Z. Cui, M.~Demirer, S.~Jaffe, L.~Musolff, S.~Peng, and T.~Salz. The effects of generative AI on
high-skilled work: Evidence from three field experiments with software developers. SSRN 4945566,
2024.

\bibitem[Bhati(2026)]{bhati2026}
H.~Bhati. Beyond code generation: Reliability, verification, and cost economics in the agentic
software development lifecycle. arXiv:2609.04681, 2026.

\bibitem[Crescitelli et al.(2026)]{crescitelli2026}
V.~Crescitelli, G.~Immediato, F.~Persia, and S.~Costantini. AI evaluation should measure
verification cost, not correctness alone. arXiv:2608.08709, 2026.

\bibitem[Koch(2026)]{agilev2026}
C.~Koch. Agentic Agile-V: From vibe coding to verified engineering in software and hardware
development. arXiv:2605.20456, 2026.

\bibitem[Juran and Godfrey(1999)]{juran1999}
J.~M. Juran and A.~B. Godfrey, editors. \emph{Juran's Quality Handbook}. McGraw-Hill, 5th
edition, 1999.

\bibitem[Feigenbaum(1961)]{feigenbaum1961}
A.~V. Feigenbaum. \emph{Total Quality Control}. McGraw-Hill, 1961.

\bibitem[Townsend(1979)]{townsend1979}
R.~M. Townsend. Optimal contracts and competitive markets with costly state verification.
\emph{Journal of Economic Theory}, 21(2):265--293, 1979.

\bibitem[Abdolmaleki et al.(2026)]{abdolmaleki2026}
M.~Abdolmaleki, S.~Jasin, and B.~Wang. Designing agentic AI workflow portfolios under imperfect
selection and compute cost. arXiv:2609.18126, 2026.

\bibitem[Demirer et al.(2026)]{demirer2026}
M.~Demirer, L.~Musolff, and L.~Yang. Writing code vs.\ shipping code: Productivity effects across
generations of AI coding tools. \emph{NBER Working Paper 35275}, doi:10.3386/w35275, 2026.

\bibitem[Becker et al.(2025)]{metr2025}
J.~Becker, N.~Rush, B.~Barnes, and D.~Rein. Measuring the impact of early-2025 AI on experienced
open-source developer productivity. arXiv:2507.09089, 2025.

\bibitem[DORA(2025)]{dora2025}
Google Cloud DORA. \emph{State of AI-Assisted Software Development 2025}. 2025.

\bibitem[Vogels(2025)]{vogels2025}
W.~Vogels. The renaissance developer. Closing keynote, AWS re:Invent, December 2025.

\bibitem[SonarSource(2026)]{sonar2026}
SonarSource. \emph{State of Code Developer Survey 2026}. SonarSource, January 2026.

\bibitem[Catalini et al.(2026)]{catalini2026}
C.~Catalini, X.~Hui, and J.~Wu. Some simple economics of AGI. arXiv:2602.20946, 2026.

\bibitem[Wang(2026)]{wang2026onebrain}
B.~Wang. One brain, many workers: Common-cause failure and the limits of redundancy in agent
fleets. Companion paper, 2026.

\bibitem[Huang et al.(2023)]{huang2023}
J.~Huang, X.~Chen, S.~Mishra, H.~Zheng, A.~W.~Yu, X.~Song, and D.~Zhou. Large language models
cannot self-correct reasoning yet. arXiv:2310.01798, 2023.

\end{thebibliography}

\end{document}
